% TOP_PROBLEM: {"id": 30006982, "title": "Zagier's dimension conjecture for classical multiple zeta values", "category_id": 1, "category": {"id": 1, "name": "number_theory", "display_name": "Number Theory", "description": "Properties of integers, prime numbers, Diophantine equations.", "slug": "number-theory", "order_index": 1, "created_at": "2026-07-31T15:26:25.670Z"}, "status": "open"}
% FIRST_POSTED: 2026-10-01
% ATTEMPT_STATUS: unresolved
% !TeX program = lualatex
% Definition and sources only; this file is not an LLM solution attempt.
\documentclass[11pt]{article}
\usepackage[margin=25mm]{geometry}
\usepackage{fontspec}
\setmainfont{Latin Modern Roman}
\usepackage{amsmath,amssymb,mathtools}
\usepackage{unicode-math}
\setmathfont{Latin Modern Math}
\usepackage{xurl,hyperref}
\hypersetup{colorlinks=true,urlcolor=blue}
\setcounter{secnumdepth}{0}
\setlength{\parindent}{0pt}
\setlength{\parskip}{5pt}
\setlength{\emergencystretch}{3em}
\begin{document}
\section*{\texorpdfstring{Zagier's dimension conjecture for classical multiple zeta values}{Zagier's dimension conjecture for classical multiple zeta values}}\label{30006982}
\subsection{Definitions and mathematical statement}
For $k_1\ge2$ and $k_2,\ldots,k_r\ge1$, set
\[
 \zeta(k_1,\ldots,k_r)=\sum_{n_1>\cdots>n_r>0}
 n_1^{-k_1}\cdots n_r^{-k_r}.
\]
Its weight is $k_1+\cdots+k_r$. Let $Z_k$ be the ${\mathbb{Q}}$-span of these real numbers at weight $k$, with $Z_0={\mathbb{Q}}$ and $Z_1=0$. Define $d_0=1,d_1=0,d_2=1$ and
$d_k=d_{k-2}+d_{k-3}$ for $k\ge3$. The target is $\dim_{{\mathbb{Q}}}Z_k=d_k$ for every $k$. These are actual numerical multiple zeta values, not formal symbols modulo chosen relations or motivic lifts; relations between different weights are not asserted absent by this dimension statement alone.
\par\smallskip\textit{Formulation and concept references:} \href{https://www2.math.kyushu-u.ac.jp/~mkaneko/papers/40conjecture_dimension.pdf}{[S1]}.

\subsection{Short English statement}
At every weight, does the vector space of actual multiple zeta values have exactly the dimension given by Zagier's two-step recurrence?
\subsection{Research attempt: the spanning recurrence and exact low-weight numerical identities}
\textbf{Status and external theorem.} We prove the combinatorial recurrence behind the proposed bound and the exact dimensions in weights two and three. For the general upper bound we use Brown's proven spanning theorem [R1] for actual multiple zeta values by arguments in $\{2,3\}$. The passage from motivic independence to independence of numerical periods is not assumed; it is the central gap in this attempt.

\subsubsection{1. Counting the proposed spanning family}
Let $a_k$ count finite words over symbols $2,3$ whose entries sum to $k$, including the empty word at weight zero. A nonempty word ends in exactly one of these symbols, so removing that symbol gives
\[a_0=1,\quad a_1=0,\quad a_2=1,\qquad a_k=a_{k-2}+a_{k-3}.\]
Consequently $a_k=d_k$, and the ordinary generating function is
\[\sum_{k\ge0}a_kt^k=\sum_{r\ge0}(t^2+t^3)^r=\frac1{1-t^2-t^3}.\]
Each such word is admissible for the numerical series because its first entry is at least two. Brown's external spanning theorem therefore yields the rigorous inequality $\dim_{\mathbb Q}Z_k\le d_k$. Counting the words alone would not establish spanning; that is precisely where the theorem is used.

\subsubsection{2. A direct proof of the weight-three reduction}
The only admissible indices at weight three are $(3)$ and $(2,1)$. Write
\[T=\sum_{m,k\ge1}\frac1{m(m+k)^2}=\zeta(2,1).\]
This converges, since its equivalent sum is $\sum_{n\ge2}H_{n-1}/n^2$ and $H_n\le1+\log n$. By interchanging $m,k$ in a second copy and adding, with all terms nonnegative,
\[2T=\sum_{m,k\ge1}\frac1{mk(m+k)}=:S.\]
For fixed $k$, partial fractions and a telescoping limit give
\[\sum_{m\ge1}\frac1{m(m+k)}=\frac{H_k}{k}.\]
Therefore
\[S=\sum_{k\ge1}\frac{H_k}{k^2}
=\sum_{k\ge1}\frac{H_{k-1}}{k^2}+\sum_{k\ge1}\frac1{k^3}
=T+\zeta(3).\]
It follows that $\zeta(2,1)=\zeta(3)$. Since $\zeta(3)>0$, the actual numerical span has dimension one, equal to $d_3$. At weight two the sole index is $(2)$ and $\zeta(2)>0$, so the same conclusion holds. We use positivity only to prove that a one-element generating family is nonzero; it supplies no general linear independence theorem.

\subsubsection{3. Why formal and motivic bases do not give the missing inequality}
Consider the vector space with basis the weight-$k$ words in $2,3$, and the map sending each word to its numerical multiple zeta value. The external spanning theorem says this evaluation map is surjective onto $Z_k$. The conjectured equality says its kernel is zero. Relations proved in an auxiliary motivic algebra may show that corresponding motivic elements are independent there, but the period homomorphism from that algebra to real numbers could, absent a further theorem, have a kernel. Independence before applying a homomorphism never implies independence after applying it.

The first small example with two proposed generators occurs at weight five, with words $(2,3)$ and $(3,2)$. Their positivity does not rule out a nonzero rational relation with opposite signs. Likewise, decimal approximations or failure to find a bounded-height relation do not prove independence over all of $\mathbb Q$. The numerical target must be distinguished from a formal shuffle/stuffle quotient whose dimension may be computable by exact linear algebra.

\subsubsection{4. Audit and remaining gap}
The double-series operations above use absolute convergence or nonnegative summands, and the harmonic-series difference is taken as a telescoping limit rather than subtraction of divergent totals. The recurrence counts ordered words, not unordered partitions of the weight. The diagnostic enumerates these words through weight thirty and checks the recurrence. The proven conclusions are the externally supported upper bound at every weight and direct equality at weights zero through three under the stipulated $Z_0,Z_1$ conventions. No general lower bound $\dim Z_k\ge d_k$ is established. Proving injectivity of numerical evaluation on the proposed spanning family remains the exact missing step.

\subsection{Sources}
\begin{itemize}
\item[S1]Masanobu Kaneko, Masayuki Noro and Ken'ichi Tsurumaki, On a conjecture for the dimension of the space of the multiple zeta values; author-hosted version, venue/DOI not verified.
\url{https://www2.math.kyushu-u.ac.jp/~mkaneko/papers/40conjecture_dimension.pdf}.
\textit{Formulation source.}
\item[R1] Francis Brown, Mixed Tate motives over Z (primary paper establishing the Hoffman spanning theorem).
\url{https://arxiv.org/abs/1102.1312}.
\end{itemize}
\end{document}
